\bookchapter{第六章：守護你的數字生活——安全與隱私}

\chapterintro{手機安全=隱私安全+財產安全。在這個資料即資產的時代，保護手機就是保護你自己。}

\booksection{一、鎖屏：第一道物理防線}

鎖屏是你手機的第一道防線，也是最容易被忽視的。

\booksubsection{安全等級排序：}

\begin{enumerate}
  \item \textbf{密碼/PIN碼}：安全度最高，是最終保障（忘記指紋/臉時靠它）。
  \item \textbf{指紋識別}：安全與便利的完美平衡。
  \item \textbf{人臉識別}：最方便，但安全性取決於技術（3D結構光安全，2D照片易騙過）。
  \item \textbf{圖案解鎖}：方便但容易留下痕跡被猜到。
\end{enumerate}

\booksubsection{設定強鎖屏：}

\begin{itemize}
  \item 務必設定一個複雜密碼（不要用生日、123456）。
  \item 開啟生物識別。
  \item 鎖屏時間設成"立即"或最短。
\end{itemize}

\booksection{二、“查詢我的手機”：終極防盜工具}

\textbf{必須開啟！}它能：

\begin{itemize}
  \item \textbf{定位}：手機丟哪了，在地圖上找。
  \item \textbf{播放聲音}：讓手機響鈴，適合在家找。
  \item \textbf{丟失模式}：遠端鎖死手機並顯示聯絡號碼。
  \item \textbf{遠端擦除}：最後一招，清除所有資料（但擦除後無法再定位）。
\end{itemize}

去 [設定] 裡，蘋果找“查詢我的iPhone”，安卓找“查詢我的裝置”或品牌類似服務，確保它是綠色開啟狀態。

\booksection{三、系統更新：打補丁，堵漏洞}

系統更新首要目的是修復安全漏洞，就像給房子補牆洞。所以，看到系統更新提示，儘量第一時間升級。

當然，大版本更新（比如iOS 17到iOS 18）可能會有相容性問題，可以觀望幾天再升。

\booksection{四、識別與防範電信詐騙}

在AI技術迅猛發展的今天，傳統的“眼見為實”認知模式已完全失效。詐騙分子正利用深度偽造（Deepfake）等AI技術，實現了從“廣撒網”式詐騙到“精準圍獵”式犯罪的全面升級。

\booksubsection{AI 時代最常見的幾種詐騙套路}

\textbf{(1) AI 換臉擬聲詐騙}\par

不法分子透過非法蒐集目標物件在社交媒體釋出的高畫質照片或僅需幾秒鐘的語音樣本，利用AI深度偽造技術克隆出足以亂真的“親友影片”或“語音”。常見劇本包括：\textbf{突發車禍、緊急借款、遭遇綁架}等。此類詐騙利用親情、友情等情感紐帶，\textbf{迷惑性極高}。

\textbf{(2) 冒充權威 / 客服詐騙}\par

AI技術能夠精準模擬官方客服、公檢法機關人員的聲音與語調，甚至可以合成高度逼真的虛假轉賬截圖、電子公文、通緝令等檔案，透過\textbf{製造恐慌心理或權威壓力}，誘導受害者轉賬或開啟螢幕共享功能，從而竊取賬戶資訊。

\textbf{(3) AI 投資理財詐騙}\par

透過AI演算法偽造“穩賺不賠”的虛假收益截圖、虛擬量化選股軟體介面或“大師級”投資分析報告，以高額回報為誘餌，售賣高價投資服務或誘導受害者將資金轉入其控制的交易賬戶，\textbf{最終捲款跑路}。

\booksubsection{破局之法——“照妖鏡”防護思路}

AI可以克隆聲音和麵容，但也不是無懈可擊，有幾個破綻可以利用：

\begin{itemize}
  \item \textbf{動作驗證法}：進行視訊通話時，主動要求對方完成預設的連貫動作指令，例如"用手在臉前快速揮過"、"側臉展示耳朵輪廓"或"連續快速眨眼5次"。目前的AI換臉技術在處理複雜動態遮擋、快速運動場景和光影變化時，\textbf{極易出現畫面卡頓、面部扭曲或邊緣穿模}等破綻。
  \item \textbf{暗號核驗法}：建立並使用預設的私密核驗問題，例如"我們上次一起聚餐時你點的什麼菜"、"我送你的第一個禮物是什麼"等只有雙方知曉的細節。若對方出現明顯遲疑或答非所問，應\textbf{立即終止通話並啟動安全預案}。
  \item \textbf{二次核實法}：嚴格遵守公安部網安局釋出的"\textbf{四不一核實}"原則——不輕信陌生資訊、不隨意轉賬匯款、不洩露個人敏感資訊、不傳播可疑內容。涉及任何錢財往來時，務必透過\textbf{已知的原號碼回撥或當面確認}。進行轉賬操作時，建議選擇"24小時到賬"模式，為潛在詐騙識別和資金挽回預留充足時間視窗。
\end{itemize}

\booksubsection{傳統常見套路（仍要警惕）：}

\begin{itemize}
  \item \textbf{中獎/刷單返利}：利誘，先給你點甜頭再騙大錢。
  \item \textbf{網戀詐騙}：殺豬盤，培養感情後騙錢。
  \item \textbf{冒充熟人}：盜號後借錢（AI 時代更難分辨）。
\end{itemize}

\booksubsection{手機自帶防護：}

\begin{itemize}
  \item 騷擾電話攔截。
  \item 驗證碼防竊取。
  \item 簡訊連結安全檢測。
\end{itemize}

\booksubsection{防騙口訣：}

\textbf{不輕信、不透露、不點選、不轉賬、有疑問先核實！}遇到“急事”先冷靜5分鐘，幾乎所有詐騙都會在這一刻露出馬腳。

\booksection{五、抵禦“開盒”與社工攻擊：斬斷隱私洩露鏈條}

“開盒”（社會工程學攻擊）的核心原理在於\textbf{系統性蒐集、分析並拼湊個人數字碎片，構建完整的個人畫像}。防範的關鍵在於為個人隱私劃定明確紅線，建立多層次防護體系。

\booksubsection{照片脫敏處理}

基礎 EXIF 清理方法見第二章“圖片與檔案系統管理”。在社交媒體分享照片前，還需進一步\textbf{徹底刪除照片的後設資料（EXIF資訊）}，包括拍攝時間、裝置型號、GPS位置座標等敏感資料。對照片中包含的\textbf{門牌號碼、車牌資訊、車票條形碼、身份證號}等關鍵資訊進行嚴格打碼或模糊處理，避免無意中洩露個人行蹤軌跡。

\begin{itemize}
  \item iOS：[設定] > [隱私與安全性] > [相機] 中可關閉"精確位置"；分享時選"僅當前照片"可剝離後設資料。
  \item 安卓：部分相簿App（如Google相簿）分享時會提示是否保留位置資訊。
  \item 通用工具：ExifTool、Metapho 等可檢視/編輯 EXIF。
\end{itemize}

\booksubsection{應用許可權收斂}

基礎許可權管理方法詳見第四章“許可權管理”。在“防開盒”場景下，還需進一步定期審查並關閉電商類、工具類App的\textbf{麥克風、通訊錄、位置資訊}等非必要許可權，防止App在後臺“偷聽”日常對話或非法收集個人關係網路；對長期閒置但仍佔用高許可權的應用程式應及時解除安裝，從源頭減少隱私洩露風險點。

\begin{itemize}
  \item 看到"通訊錄許可權"申請時，先問一句：這App真有必要知道我所有聯絡人嗎？
  \item 三個月沒用過的App，果斷解除安裝。許可權是越用越緊的，不是越用越松。
\end{itemize}

\booksubsection{智慧裝置安全設防}

智慧攝像頭、智慧門鎖、智慧音箱等IoT裝置在初次使用時必須\textbf{立即修改預設登入密碼}，設定包含大小寫字母、數字和特殊符號的高強度獨立密碼；不常用的遠端訪問功能應保持關閉狀態，定期檢查裝置韌體更新並及時修補安全漏洞，\textbf{防止駭客入侵併窺探家庭隱私空間}。

\begin{tcbwarn}{真實案例警示}
2022年某地警方通報，駭客透過破解某品牌家用攝像頭的預設密碼（很多使用者從未修改），竊取了上千戶家庭的臥室畫面並在網上兜售。\textbf{不改預設密碼 = 把家大門敞開}。

\end{tcbwarn}

\booksection{六、隱私保護：你的資料是你的財產}

手機裡的照片、聊天記錄、位置軌跡，本質上都是你的“數字資產”。一旦洩露，輕則廣告騷擾，重則被人肉、詐騙甚至敲詐。

\booksubsection{App 許可權的“最小化原則”：}

\begin{itemize}
  \item \textbf{位置}：打車/地圖用"始終允許"；其他 App 用"僅使用時允許"。
  \item \textbf{通訊錄}：除非必要（如微信加好友）才給，否則一律禁止。
  \item \textbf{相簿 / 全部照片}：iOS 14+ 起可授權"選定的照片"，安卓 13+ 同樣支援精細授權。
  \item \textbf{麥克風 / 攝像頭}：平時關掉，需要時再開啟。懷疑被偷錄時，滑出控制中心看是否有"小綠點/小橙點"。
  \item \textbf{後臺跟蹤}：iOS 14.5+ 起強制要求 App 使用 IDFA 前必須詢問；安卓可關閉"廣告 ID"。
\end{itemize}

\booksubsection{隱私管理實戰：}

\begin{itemize}
  \item \textbf{隱私報告（iOS）}：[設定] > [隱私與安全性] > [App 隱私報告]，檢視哪些 App 在偷偷訪問你的位置、相機、麥克風。
  \item \textbf{許可權使用記錄（安卓）}：[設定] > [隱私] > [許可權使用記錄]，看誰在什麼時候呼叫了敏感許可權。
  \item \textbf{限制廣告跟蹤}：在系統設定裡關閉"廣告 ID"或"限制廣告追蹤"，能減少"被精準廣告騷擾"。
\end{itemize}

\booksection{七、丟失 / 被盜後的“急救”清單}

手機丟了不要慌，按以下順序處理：

\begin{enumerate}
  \item \textbf{立即凍結 SIM 卡}：聯絡運營商掛失，避免被用於接收驗證碼。
  \item \textbf{遠端鎖定}：透過雲服務（iCloud / 品牌雲）開啟"丟失模式"。
  \item \textbf{修改重要密碼}：登入過手機的所有重要賬號（支付、銀行、郵箱）立即改密碼。
  \item \textbf{報案}：到派出所報案，記錄 IMEI 號碼。手機找回時也便於證明身份。
  \item \textbf{通知親友}：防止騙子以你名義借錢。
  \item \textbf{（最後）遠端擦除}：確認找回無望時再操作，擦除後定位功能失效。
\end{enumerate}

\booksection{八、SIM 卡與 eSIM 安全}

\begin{itemize}
  \item \textbf{設定 SIM 卡 PIN}：開啟後，每次重啟手機或換卡都需要輸入 PIN 碼，防止"拔卡換手機"繞過鎖屏。
  \item \textbf{eSIM 優勢}：沒有實體卡可被拔出，安全性更高；同時支援多號切換，旅行時非常方便。
  \item \textbf{二次換卡風險}：警惕"運營商二次換卡"騙局——騙子冒用你的身份補辦 SIM 卡，從而接管你的手機號。設定"補卡需機主本人持身份證到營業廳辦理"是最有效的防線。
\end{itemize}

\begin{tcbwarn}{千萬別做的事}
\begin{itemize}
  \item 用生日 / 123456 / qwerty 等弱密碼；
  \item 把驗證碼發給任何人（哪怕是"客服"）；
  \item 點選簡訊裡的不明連結，特別是 00 開頭的境外號碼；
  \item 在公共 Wi-Fi 下登入網銀或支付；
  \item 給陌生人開啟"螢幕共享"或"遠端控制"。
\end{itemize}

\end{tcbwarn}

\begin{tcbpractice}{試著啟用並模擬"遠端鎖定"}
\begin{enumerate}
  \item 確保你的"查詢我的手機"功能已開啟。
  \item 在電腦瀏覽器或另一部手機上，登入對應的雲服務賬號（iCloud.com 或品牌雲服務網站）。
  \item 找到"查詢裝置"功能，嘗試對你自己的手機發起一次"播放聲音"或模擬"鎖定"指令，熟悉整個操作流程。
\end{enumerate}

\end{tcbpractice}
